% 附录
\chapter{关键实现说明}

由于完整工程代码体量较大，且包含训练脚本、模型定义、编码调用与评估工具等多个部分，为避免影响论文主体的连贯性，本文仅在附录中给出核心模块的详细结构、伪代码、主要参数配置及关键运行命令。完整代码作为毕业设计附件单独提交。

% ============================================================
\section{RAPNet (Actor) 详细网络结构}
% ============================================================

RAPNet采用``Pixel Unshuffle $\rightarrow$ 残差卷积 $\rightarrow$ Pixel Shuffle''的三段式架构，在低分辨率特征空间完成残差增强后映射回原分辨率。各层的具体参数如表~\ref{tab:rapnet-detail}所示。

\begin{table}[H]
  \centering
  \small
  \setlength{\tabcolsep}{4pt}
  \caption{RAPNet (Actor) 网络结构详情（$\text{scale}=2$时，$C_{in}=3 \times 2^2 = 12$）}
  \label{tab:rapnet-detail}
  \resizebox{\textwidth}{!}{%
  \begin{tabular}{llccl}
    \toprule
    \textbf{阶段} & \textbf{层名称} & \textbf{输入$\rightarrow$输出通道} & \textbf{核/步长/填充} & \textbf{激活} \\
    \midrule
    色度上采样 & \texttt{upsample\_chroma} & $(1,H/2,W/2)\rightarrow(1,H,W)$ & Bilinear $\times 2$ & — \\
    \midrule
    空间折叠 & \texttt{PixelUnshuffle(2)} & $(3,H,W)\rightarrow(12,H/2,W/2)$ & — & — \\
    \midrule
    主干入口 & \texttt{head} (Conv2d) & $12 \rightarrow 16$ & $3\times3$ / 1 / 1 & LeakyReLU(0.1) \\
    \midrule
    \multirow{2}{*}{残差体$\times 5$} & \texttt{body[i].conv1} & $16 \rightarrow 16$ & $3\times3$ / 1 / 1 & LeakyReLU(0.1) \\
     & \texttt{body[i].conv2} & $16 \rightarrow 16$ & $3\times3$ / 1 / 1 & 残差连接 \\
    \midrule
    主干出口 & \texttt{tail} (Conv2d) & $16 \rightarrow 12$ & $3\times3$ / 1 / 1 & — \\
    \midrule
    空间展开 & \texttt{PixelShuffle(2)} & $(12,H/2,W/2)\rightarrow(3,H,W)$ & — & — \\
    \midrule
    残差限幅 & $\delta = \tanh(\cdot) \times 0.05$ & — & — & Tanh \\
    色度均值归零 & $\delta_{Cb/Cr} \mathrel{-}= \text{mean}(\delta_{Cb/Cr})$ & — & — & — \\
    \midrule
    色度下采样 & \texttt{downsample\_chroma} & $(1,H,W)\rightarrow(1,H/2,W/2)$ & Bilinear $\times 0.5$ & — \\
    \bottomrule
  \end{tabular}
  }
\end{table}

总参数量约113K。其中5个ResBlock贡献了主要参数量（每个ResBlock含两层$16\rightarrow16$的Conv2d，参数$2\times(16\times16\times3\times3+16) = 4\,640$，五个共$23\,200$），head/tail的$12\leftrightarrow16$通道转换各贡献约$1\,744 / 1\,740$参数。

% ============================================================
\section{ValueNet (Critic) 详细网络结构}
% ============================================================

ValueNet采用双分支对称结构，分别对原始帧（状态$s$）和增强帧（动作$a$）提取特征，拼接后进行融合，最终输出标量Q值。各层的具体参数如表~\ref{tab:valuenet-detail}所示。

\begin{table}[!htbp]
  \centering
  \small
  \setlength{\tabcolsep}{4pt}
  \caption{ValueNet (Critic) 网络结构详情}
  \label{tab:valuenet-detail}
  \resizebox{\textwidth}{!}{%
  \begin{tabular}{llccl}
    \toprule
    \textbf{阶段} & \textbf{层名称} & \textbf{输入$\rightarrow$输出通道} & \textbf{核/步长/填充} & \textbf{激活} \\
    \midrule
    \multirow{2}{*}{输入分支} & \texttt{branch\_in.conv1} & $3 \rightarrow 32$ & $3\times3$ / 1 / 1 & LeakyReLU(0.1) \\
     & \texttt{branch\_in.conv2} & $32 \rightarrow 32$ & $3\times3$ / 1 / 1 & — \\
    \midrule
    \multirow{2}{*}{输出分支} & \texttt{branch\_out.conv1} & $3 \rightarrow 32$ & $3\times3$ / 1 / 1 & LeakyReLU(0.1) \\
     & \texttt{branch\_out.conv2} & $32 \rightarrow 32$ & $3\times3$ / 1 / 1 & — \\
    \midrule
    特征拼接 & \texttt{cat(in, out)} & $32+32=64$ & — & — \\
    融合卷积 & \texttt{fusion\_conv} & $64 \rightarrow 64$ & $3\times3$ / 1 / 1 & LeakyReLU(0.1) \\
    \midrule
    全局池化 & \texttt{AdaptiveAvgPool2d(1)} & $(64,H,W)\rightarrow(64,1,1)$ & — & — \\
    时间平均 & $\text{mean}(\cdot, \text{dim}=T)$ & $(B,T,64)\rightarrow(B,64)$ & — & — \\
    \midrule
    \multirow{2}{*}{MLP头} & \texttt{head.fc1} & $64 \rightarrow 64$ & — & ReLU \\
     & \texttt{head.fc2} & $64 \rightarrow 1$ & — & — \\
    \bottomrule
  \end{tabular}
  }
\end{table}

% ============================================================
\Needspace{0.7\textheight}
\section{损失权重与训练超参数}
% ============================================================

表~\ref{tab:hyperparams}汇总了模型在DDPG风格Actor-Critic训练框架中的完整超参数配置。这些参数来自最终收敛版本（v44），此前经历了近50个版本的迭代调参。

\begingroup
\small
\setlength{\tabcolsep}{4pt}
\setlength{\LTleft}{0pt plus 1fill}
\setlength{\LTright}{0pt plus 1fill}
\begin{longtable}{lll}
  \caption{训练主要超参数配置}
  \label{tab:hyperparams}\\
  \toprule
  \textbf{类别} & \textbf{超参数} & \textbf{设定值} \\
  \midrule
  \endfirsthead
  \caption[]{训练主要超参数配置（续）}\\
  \toprule
  \textbf{类别} & \textbf{超参数} & \textbf{设定值} \\
  \midrule
  \endhead
  \midrule
  \multicolumn{3}{r}{续下页} \\
  \endfoot
  \bottomrule
  \endlastfoot
  \multirow{4}{*}{迭代与数据}
    & 训练总迭代次数 & $20\,000$ \\
    & 每次采样帧数 (clip\_len) & 16 \\
    & 输入分辨率 & $1920 \times 1080$ \\
    & Random Crop裁剪尺寸 & $960 \times 540$ \\
  \midrule
  \multirow{4}{*}{学习率}
    & Actor学习率 $\eta_{Actor}$ & $1 \times 10^{-5}$ \\
    & Critic学习率 $\eta_{Critic}$ & $2 \times 10^{-5}$ \\
    & 学习率衰减步长 & 每1000 iter \\
    & 学习率衰减因子 $\gamma_{lr}$ & $1/3$ \\
  \midrule
  \multirow{3}{*}{奖励函数}
    & 质量权重 $\alpha$ & $200$ \\
    & 码率权重 $\beta$ & $100$ \\
    & QP测试点集 $\mathcal{Q}$ & $\{37, 40, 42, 45, 47\}$ \\
  \midrule
  \multirow{3}{*}{Actor损失权重}
    & MSE结构损失 $w_{mse}$ & $5$ \\
    & 非边缘高频抑制 $w_{hf}$ & $10$ \\
    & TV平滑损失 $w_{TV}$ & $0.1$ \\
  \midrule
  \multirow{2}{*}{残差限幅}
    & Tanh缩放系数 & $0.05$（即$\pm 5\%$） \\
    & Cb/Cr通道 & 零均值约束 \\
  \midrule
  \multirow{2}{*}{硬件分配}
    & Actor & GPU 0 (cuda:0) \\
    & Critic & GPU 1 (cuda:1) \\
\end{longtable}
\endgroup

\FloatBarrier

% ============================================================
\section{FFmpeg编码/解码命令示例}
% ============================================================

训练环路中的编码与解码通过Python \texttt{subprocess}调用FFmpeg完成。以下是从代码中提取的典型FFmpeg命令结构：

\noindent\textbf{编码命令}（NVENC HEVC硬件编码，恒定QP模式）：
\begin{lstlisting}[language=bash,basicstyle=\ttfamily\small,breaklines=true]
ffmpeg -y \
  -f rawvideo -pix_fmt yuv420p \
  -s:v 960x540 -r ${fps} \
  -i input_crop.yuv \
  -c:v hevc_nvenc \
  -rc constqp -qp 42 \
  -preset p4 -profile:v main \
  -pix_fmt yuv420p \
  output.hevc
\end{lstlisting}

\noindent\textbf{解码命令}（HEVC $\rightarrow$ 原始YUV420）：
\begin{lstlisting}[language=bash,basicstyle=\ttfamily\small,breaklines=true]
ffmpeg -y \
  -i output.hevc \
  -f rawvideo -pix_fmt yuv420p \
  decoded.yuv
\end{lstlisting}

\noindent 其中，\texttt{-rc constqp}表示恒定QP码率控制模式。训练阶段遍历5个QP点，即$\mathcal{Q}=\{37,40,42,45,47\}$，并对奖励取平均。\texttt{-preset p4}对应NVENC的Medium编码预设，\texttt{-r \${fps}}表示按当前序列原始帧率设置时间基准，例如BasketballDrive/Cactus为50fps，ParkScene/Kimono1为24fps，BQTerrace为60fps。

% ============================================================
\section{核心代码文件目录说明}
% ============================================================

随论文提交的代码附件已整理为\texttt{毕业论文代码/}目录。该目录只包含论文复现实验所需的源码、工具脚本、依赖说明与模型权重；标准YUV测试序列体积较大，需按本地环境另行放置并修改脚本中的数据集路径。主要文件组织结构如下：

\begin{lstlisting}[basicstyle=\ttfamily\small]
毕业论文代码/
+-- README.md
+-- requirements.txt
+-- src/
|   +-- rapnet_hevc_ppo_qp_yuv_1080p.py
|   |     # 最终1080p Random Crop训练脚本
|   |     # 包含: RAPNetYCbCr (Actor),
|   |     #       ValueNetYCbCrPair (Critic),
|   |     #       DeterministicACAgentYCbCr,
|   |     #       NvencClipEnvQP,
|   |     #       train_rapnet_with_config()
|   +-- rapnet_hevc_ppo_qp_yuv_simple.py
|   |     # 下采样版本训练与完整视频评估依赖脚本
|   +-- gradient_monitor.py
|         # Actor-Critic梯度诊断工具
+-- tools/
|   +-- eval_rapnet_full_video.py
|   |     # 完整视频推理与RD评估脚本
|   +-- evaluate_time.py
|   |     # 读入、前处理、编码、解码端到端耗时统计
|   +-- run_ablation_3variants.py
|         # 三组损失配置消融实验自动化脚本
+-- checkpoints/
|   +-- actor_best.pth
|   +-- critic_best.pth
+-- results/
|     # 放置Baseline RD曲线JSON与评估输出
+-- tmp/
      # FFmpeg编码/解码临时文件目录
\end{lstlisting}

\FloatBarrier

% ============================================================
\chapter{代码提交说明}
% ============================================================

本课题完整代码已随毕业设计材料以\texttt{毕业论文代码/}目录形式提交。代码附件主要包括以下文件：

\begin{itemize}
  \item \path{src/rapnet_hevc_ppo_qp_yuv_1080p.py}：最终采用的1080p Random Crop训练入口；
  \item \path{src/rapnet_hevc_ppo_qp_yuv_simple.py}：下采样版本训练与完整视频评估所需的公共函数；
  \item \path{tools/eval_rapnet_full_video.py}：完整视频推理与RD评估脚本；
  \item \path{tools/evaluate_time.py}：端到端耗时统计脚本；
  \item \path{tools/run_ablation_3variants.py}：三组损失配置消融实验脚本；
  \item \path{checkpoints/actor_best.pth}与\path{checkpoints/critic_best.pth}：最终提交的最优Actor与Critic权重。
\end{itemize}

为避免论文附录过长，本文不在附录中逐行粘贴完整Python源码，而是给出关键网络结构、训练参数、编码命令和文件组织说明。复现实验时，可先根据\path{requirements.txt}安装依赖，再准备HEVC CTC Class B和UVG等YUV420测试序列，并将脚本中的\texttt{VIDEO\_META\_LIST}路径改为本机数据集路径。Baseline RD曲线JSON可由训练脚本中的Baseline构建函数生成，也可放置于\path{results/}目录后通过环境变量\texttt{RAPNET\_BASELINE\_JSON}指定。

如需复现实验结果，请确保以下环境配置：
\begin{itemize}
  \item NVIDIA GPU $\times$ 2（推荐RTX 4090或同等级别，显存$\geq$16GB）
  \item CUDA 12.x + cuDNN 9.x
  \item Python 3.10+, PyTorch 2.0+
  \item FFmpeg 6.0+（需包含\texttt{hevc\_nvenc}编码器支持）
  \item 关键Python包：\texttt{lpips}, \texttt{tensorboard}, \texttt{numpy}, \texttt{opencv-python}
\end{itemize}
